% try: add a \section, or an \index{...} entry: the contents and the index follow in one build
\documentclass[11pt]{book}
\usepackage[paperwidth=6in,paperheight=9in,inner=.7in,outer=1.45in,top=.8in,bottom=.8in,marginparwidth=1in,marginparsep=.15in]{geometry}
\usepackage{lettrine}
\usepackage{makeidx}
\makeindex

\newcommand{\epigraph}[2]{\begin{flushright}\begin{minipage}{.7\textwidth}\itshape\small #1\par\vspace{2pt}\upshape\hfill--- #2\end{minipage}\end{flushright}\vspace{1.5em}}

\begin{document}
\frontmatter
\tableofcontents

\mainmatter
\chapter{The Shape of a Page}
\epigraph{Beauty is our business.}{Donald E. Knuth}

\lettrine[lines=3]{T}{he page} is the oldest interface we have. Long before
screens, typographers\index{typography} learned how the eye moves down a
column, how far apart lines should sit, and where a word may break. A
typesetter\index{typesetter} keeps all of that in its rules, so the writer
does not have to.\marginpar{\footnotesize Try editing this paragraph: only it is set again.}

\section{Lines}
TeX breaks a paragraph into lines by looking at the whole paragraph at
once\index{line breaking}, weighing every possible break against the
others. The result is even spacing with few hyphens\index{hyphenation}, at
the cost of reading the paragraph to its end before setting its first line.

\section{Pages}
Pages are built\index{page building} from the lines, with floats and
footnotes\footnote{Like this one.} placed where they fit. An edit early in a
book can move every page after it; PhiTeX sets them again only until they
match what it built before.

\section{The Index}
This short chapter has an index\index{index} at the back, made by
\texttt{makeindex} inside the build.

\backmatter
\printindex
\end{document}
